% ==================== 6.1 引言 ====================
\section{引言}
\label{sec:ch6_intro}

路径规划是巡视器自主行走的核心功能之一。在数字孪生环境下，路径规划可以利用虚拟空间中的环境模型、动力学模型等资源，实现更高效、更安全的路径生成。本章研究基于数字孪生的巡视器路径规划方法。

巡视器路径规划面临以下挑战：

\textbf{（1）环境不确定性}

月面环境复杂多变，地形、月壤特性存在空间变化，障碍物分布未知。规划算法需要在不确定性条件下做出决策。

\textbf{（2）动力学约束}

巡视器的运动受动力学约束，包括最大速度、最大加速度、转弯半径等。规划路径必须满足这些约束才能实际执行。

\textbf{（3）多目标优化}

路径规划需要同时考虑路径长度、能耗、安全性、时间等多个目标，这些目标往往相互冲突，需要进行权衡。

\textbf{（4）实时性要求}

在动态环境下，规划算法需要快速响应环境变化，实时更新路径。

针对上述挑战，本章构建数字孪生环境下的路径规划框架，提出基于A-D3QN的混合路径规划算法，实现全局规划与局部规划的协同优化。
